\documentclass[11pt,a4paper]{article}
\usepackage{ctex}
\usepackage{times}
\usepackage{amsmath}
\usepackage{booktabs}
\usepackage{graphicx}
\usepackage{tikz}
\usepackage{algpseudocode}
\usepackage{algorithm}
\usepackage{xcolor}
\usepackage{hyperref}

\setlength{\topmargin}{-0.5in}
\setlength{\textheight}{9in}
\setlength{\oddsidemargin}{-0.3in}
\setlength{\textwidth}{7in}

\newcommand{\method}{PersonalQuery}
\newcommand{\did}{\text{DiD}}
\newcommand{\didscore}{\text{Score}_{\text{DiD}}}

\title{\textbf{接地气个性化搜索查询生成：基于语言风格对齐的多阶段流水线}}

\author{
匿名作者 \\
单位名称 \\
\texttt{email@institution.edu}
}

\begin{document}
\maketitle

\begin{abstract}
个性化搜索系统传统上依赖简单的用户画像或协同过滤，往往无法捕捉个体用户微妙偏好和沟通模式的细微差别。我们介绍了PersonalQuery，这是一个新颖的12阶段流水线，通过从历史评论中提取细粒度用户偏好并将查询语言与个人写作模式对齐，生成"接地气"的个性化搜索查询。我们的方法集成了基于大语言模型（LLM）的偏好提取、多维度语言风格分析（16个风格特征）、迭代式风格优化（特征感知提示）和卷积神经网络（CNN）拼写难度预测，以实现靶向噪声注入。在Amazon产品评论上的实验评估表明，与基线方法相比，PersonalQuery生成的查询在语义相关性和风格真实性方面都取得了显著改进，人类评估证实了LLM和人类判断之间存在强相关性（Spearman's $\rho = 0.82$）。

\textbf{关键词}：个性化搜索、查询生成、用户建模、语言风格迁移、大语言模型
\end{abstract}

\section{引言}

\subsection{背景与动机}

电商搜索系统充当用户与庞大产品目录之间的主要接口。传统的搜索界面一视同仁地对待所有用户，对于相同的查询返回相同的结果，无论个体偏好、专长水平或沟通模式如何。这种"一刀切"的方法无法捕捉用户表达需求和评估产品的丰富异质性。

近年来，个性化搜索的进展主要集中在基于用户历史或协作信号的重新排序上。然而，这些方法忽略了一个关键维度：查询本身。用户在搜索查询中使用的语言不仅反映了他们的信息需求，还反映了他们的沟通风格、技术娴熟程度和个人偏好。一个专业手工艺人搜索"3mm德式玻璃闪粉，80目纹理"与一个初学者询问"闪亮的工艺闪粉"的沟通方式完全不同。

\subsection{接地气个性化挑战}

我们将"接地气个性化"定义为生成满足以下条件的搜索查询：
\begin{enumerate}
    \item \textbf{语义忠实}于从行为数据中提取的用户偏好
    \item \textbf{风格真实}反映个人写作模式
    \item \textbf{行为现实}包含自然变化，如拼写模式
\end{enumerate}

实现接地气个性化需要解决几个挑战：
\begin{itemize}
    \item \textbf{细粒度偏好提取}：超越粗粒度类别，到与情感相关的特定属性
    \item \textbf{风格建模}：捕捉词汇之外的语言模式，包括句法复杂性和话语特征
    \item \textbf{可控生成}：平衡语义准确性与风格对齐
    \item \textbf{现实变化}：复制用户特定的怪癖，如拼写倾向
\end{itemize}

\subsection{贡献}

我们介绍了PersonalQuery，这是一个综合性的流水线，通过以下方式解决这些挑战：

\begin{enumerate}
    \item \textbf{12阶段处理流水线}：将原始用户评论转化为具有可控风格属性的个性化搜索查询

    \item \textbf{多维度语言特征分析}：使用16个风格特征，包括句法复杂性、词汇密度和词性分布

    \item \textbf{迭代式风格优化（特征感知提示）}：通过差距分析逐步将生成的查询与目标风格画像对齐

    \item \textbf{基于CNN的拼写难度模型}：预测单词级易错性，以实现靶向噪声注入

    \item \textbf{综合评估框架}：结合基于LLM的评估和人类对齐指标
\end{enumerate}

\section{相关工作}

\subsection{个性化搜索与推荐}

个性化搜索系统已经从简单的用户画像发展到复杂的神经方法。早期工作集中在基于用户历史的点击率预测和结果重新排序。较新的方法采用深度学习来建模用户偏好，但通常在结果级别而非查询生成级别操作。

\subsection{用户偏好提取}

从文本中提取细粒度偏好一直通过基于方面的情感分析和观点挖掘进行研究。传统方法依赖手工特征和词典，而现代方法利用预训练语言模型。我们的方法通过使用精心设计的提示工程从LLM提取结构化属性-情感对，扩展了这项工作。

\subsection{语言风格迁移与语言分析}

语言风格迁移已应用于各种领域，包括正式程度调整、情感修改和作者模仿。基于特征的方法通常依赖手工特征，而神经方法隐式地学习风格表示。我们的工作通过使用显式的16维特征配置来指导迭代式LLM优化，连接了这两种方法。

\subsection{查询生成与扩展}

查询生成已用于对话搜索和查询建议。大多数方法关注语义扩展或澄清问题。我们的工作通过生成匹配用户特定语言风格同时保持语义意图的查询来区分于现有方法。

\section{方法}

\subsection{系统概览}

PersonalQuery实现了一个12阶段流水线，通过偏好提取、画像生成、查询合成、风格优化和噪声注入来处理用户评论（图1）。每个阶段都设计为模块化且可独立评估。

\begin{figure}[h]
\centering
\begin{tikzpicture}[
    node distance=0.6cm,
    stage/.style={rectangle, draw, rounded corners, minimum width=4cm, minimum height=0.7cm, align=center, font=\small},
    arrow/.style={->, >=stealth, thick}
]
    \node[stage, fill=blue!15] (s01) {阶段0-1: 数据准备与偏好提取};
    \node[stage, fill=green!15, below=of s01] (s23) {阶段2-3: 分类与划分};
    \node[stage, fill=yellow!15, below=of s23] (s4) {阶段4: 接地气画像生成};
    \node[stage, fill=orange!15, below=of s4] (s56) {阶段5-6: 风格分析与特征提取};
    \node[stage, fill=red!15, below=of s56] (s78) {阶段7-8: 查询生成与优化};
    \node[stage, fill=purple!15, below=of s78] (s910) {阶段9-10: 噪声注入};
    \node[stage, fill=gray!15, below=of s910] (s1112) {阶段11-12: 评估};

    \draw[arrow] (s01) -- (s23);
    \draw[arrow] (s23) -- (s4);
    \draw[arrow] (s4) -- (s56);
    \draw[arrow] (s56) -- (s78);
    \draw[arrow] (s78) -- (s910);
    \draw[arrow] (s910) -- (s1112);

    \node[left=0.3cm of s01, font=\footnotesize] (input) {原始评论};
    \node[right=0.3cm of s1112, font=\footnotesize] (output) {个性化查询};
    \draw[arrow] (input) -- (s01);
    \draw[arrow] (s1112) -- (output);
\end{tikzpicture}
\caption{PersonalQuery流水线架构}
\label{fig:pipeline}
\end{figure}

\subsection{阶段0-1：数据准备与偏好提取}

\subsubsection{用户选择标准}

我们基于以下标准从Amazon评论数据中选择高质量用户：
\begin{itemize}
    \item 评论数量在[100, 110]产品范围内
    \item 产品具有有效元数据（排除"Unknown"类别）
    \item 每个产品至少5条评论，以确保可靠的偏好推断
    \item 每个产品至少4条其他用户的评论，用于公开属性提取
\end{itemize}

\subsubsection{基于LLM的偏好提取}

对于每个用户-产品对，我们采用大语言模型（GLM-5）从评论文本中提取结构化偏好。提取提示指导模型识别：

\begin{verbatim}
{
  "target_user_preferences": {
    "Use Case": [{"entity": "card making", "sentiment": "positive"}],
    "Quality": [{"entity": "adhesive strength", "sentiment": "positive"}],
    "Design": [{"entity": "color vibrancy", "sentiment": "positive"}]
  },
  "other_users_preferences": {
    "Quality": [{"entity": "durability", "sentiment": "positive"}]
  },
  "public_attributes": [
    {
      "attribute": "price range",
      "sentiment": "neutral"
    }
  ]
}
\end{verbatim}

提取保留每个属性的情感极性（积极/消极/中性），实现情感感知的查询生成。

\subsection{阶段2-3：偏好分类与数据划分}

\subsubsection{目标属性vs公开属性}

我们区分：
\begin{itemize}
    \item \textbf{目标属性}：目标用户独有的偏好
    \item \textbf{公开属性}：其他用户共有的偏好
\end{itemize}

这种区分使我们能够生成个性化查询（目标聚焦）和基线公开查询（通用偏好）。

\subsubsection{语义去重}

为了防止属性冗余，我们应用三级去重：
\begin{enumerate}
    \item \textbf{精确匹配}：四元组（asin, category, attribute, sentiment）去重
    \item \textbf{字符串相似度}：70\% Levenshtein相似度阈值
    \item \textbf{语义聚类}：句子BERT嵌入+社区检测（阈值0.85）
\end{enumerate}

\subsection{阶段4：接地气画像生成}

\subsubsection{技能水平估算}

我们基于以下标准对用户技能水平进行分类（初学者/中级/高级/专家）：
\begin{itemize}
    \item 关键词匹配（如"初学者"、"专业"）
    \item 技术术语检测（如毫米单位、克每平方米纸重）
    \item 设备/品牌提及（如"Cricut"、"Sizzix"）
\end{itemize}

\subsubsection{使用场景分类}

属性使用关键词匹配映射到预定义的使用场景：
\begin{itemize}
    \item greeting\_cards（问候卡片）
    \item scrapbooking（剪贴簿）
    \item jewelry（首饰）
    \item costume（戏服）
    \item fabric\_crafts（布料工艺）
    \item home\_decor（家居装饰）
    \item gift（礼物）
\end{itemize}

\subsubsection{情感配置}

我们计算所有偏好的情感分布：
\begin{itemize}
    \item 积极/消极/中性比例
    \item 性格特征推断（如"通常满意"、"挑剔且注重细节"）
    \item 按偏好类别的维度情感分析
\end{itemize}

\subsubsection{画像综合}

最终画像提示整合：
\begin{itemize}
    \item 产品类别
    \item 技能水平及指标
    \item 主要使用场景及示例
    \item 评价风格和性格特征
    \item 维度情感洞察
    \item 价值优先级（维度分数）
\end{itemize}

LLM生成80-120字的接地气画像描述。

\subsection{阶段5-6：风格分析与特征提取}

\subsubsection{16维语言特征}

我们使用ProfilingUD提取以下特征：

\begin{table}[h]
\centering
\small
\begin{tabular}{ll}
\toprule
\textbf{类别} & \textbf{特征} \\
\midrule
\textbf{长度} & tokens\_per\_sent, char\_per\_tok, n\_tokens \\
\textbf{词汇} & ttr\_lemma\_chunks\_100, lexical\_density \\
\textbf{词性分布} & NOUN, VERB, ADJ, ADV, PRON, DET, AUX, PART, SCONJ, CCONJ, ADP \\
\bottomrule
\end{tabular}
\caption{16维风格特征}
\label{tab:features}
\end{table}

特征归一化到[0, 1]范围用于距离计算。

\subsection{阶段7：双重查询生成}

\subsubsection{个性化查询生成}

对于个性化查询，我们：
\begin{enumerate}
    \item 选择3个具有情感平衡的属性（最多2个积极/中性+1个消极）
    \item 生成反映属性情感的第一人称查询
    \item 强制25-30字的长度约束
\end{enumerate}

提示结构：
\begin{verbatim}
任务：基于用户的属性及其情感生成个性化搜索查询。

产品类别：{category}
用户特定属性及情感：{attrs_str}

要求：
- 反映用户对每个属性的情感
- 长度必须正好是25-30字
- 使用第一人称视角
- 自然且对话式

输出格式：JSON，包含"query"和"word_count"字段
\end{verbatim}

\subsubsection{公开查询生成}

公开查询使用其他用户的3个最频繁属性，遵循相同的格式约束。

\subsection{阶段8：迭代式风格优化}

\subsubsection{特征感知提示}

优化过程使用多轮生成和靶向指令：

\textbf{第0轮}：基础生成，提供风格描述
\textbf{第1轮+}：差距驱动优化

对于每轮：
\begin{enumerate}
    \item 从当前最佳查询提取特征
    \item 计算与用户画像的特征差距
    \item 为前Top 5差距生成靶向指令
    \item 创建带具体调整的优化提示
\end{enumerate}

\subsubsection{候选选择}

对于每轮，我们使用组合分数选择最佳候选：

\begin{equation}
\text{score} = 0.7 \cdot d_{\text{style}} + 0.3 \cdot d_{\text{semantic}}
\end{equation}

其中：
\begin{itemize}
    \item $d_{\text{style}}$：查询与用户特征向量的余弦距离
    \item $d_{\text{semantic}}$：查询与原始查询的语义距离（1 - cosine\_similarity）
\end{itemize}

\subsubsection{收敛标准}

优化在以下条件下停止：
\begin{enumerate}
    \item 同时满足两个阈值：$d_{\text{style}} < 0.3$ AND $d_{\text{semantic}} < 0.4$
    \item 无改进：improvement $< 0.02$
    \item 达到最大轮数（默认：5）
\end{enumerate}

\subsection{阶段9：拼写难度模型}

\subsubsection{模型架构}

我们训练基于CNN的模型来预测单词级拼写难度：

\begin{itemize}
    \item 字符嵌入+Conv1D层（128$\rightarrow$512维）
    \item 手工特征：单词长度、音节数、频率
    \item 用户特征：错误类型频率
    \item 输出：难度分数[0, 1]
\end{itemize}

\subsection{阶段10：靶向噪声注入}

\subsubsection{错误分布建模}

对于每个用户，我们计算错误类型频率：
\begin{itemize}
    \item 同音混淆：35\%
    \item 字母遗漏：25\%
    \item 字母替换：20\%
    \item 字母插入：15\%
    \item 位置调换：5\%
\end{itemize}

\subsubsection{难度导向注入}

使用训练好的拼写模型：
\begin{enumerate}
    \item 对查询中的每个单词评分难度
    \item 选择难度最高的Top-K单词
    \item 根据用户的错误分布注入错误
    \item LLM验证语义保持
\end{enumerate}

\subsection{阶段11-12：评估}

\subsubsection{LLM 5规则评估}

我们采用基于5个标准的LLM评估：
\begin{enumerate}
    \item \textbf{偏好对齐}：查询反映用户偏好？
    \item \textbf{画像一致性}：查询匹配用户画像？
    \item \textbf{语义完整性}：涵盖所有关键属性？
    \item \textbf{自然性}：查询语言自然？
    \item \textbf{特异性}：查询适当具体？
\end{enumerate}

每个标准评分1-5，最大总分25。

\subsubsection{\textbf{基础相似度偏差修正}}

评估个性化查询时的一个关键挑战是\textbf{基础相似度偏差}：由于目标用户（$P_u$）本质上是大众的一部分，他们的画像不可避免地与大众画像（$P_g$）存在特征重叠。这导致公开查询（$Q_g$）在评估为$P_u$时天然获得"保底分"，使得难以分离个性化查询的真正增量价值。

我们引入四种互补方法来修正这种偏差：

\textbf{方法1：双重差分法}

这是最严谨的统计学修正方法。我们不看绝对分数，而是看交叉对比的相对差异。给定LLM给出的四个基础分数$S(Q, P)$：

\begin{itemize}
    \item $\Delta_u = S(Q_u, P_u) - S(Q_u, P_g)$：个性化查询的"特异性倾向"
    \item $\Delta_g = S(Q_g, P_u) - S(Q_g, P_g)$：公开查询的"特异性倾向"（通常$\approx$0）
\end{itemize}

双重差分分数：

\begin{equation}
\text{Score}_{\text{DiD}} = \Delta_u - \Delta_g
\end{equation}

原理解释：这个公式直接抵消了因为$P_u$和$P_g$相似而带来的基础得分。只有当$Q_u$捕捉到$P_u$中异于$P_g$的独特维度时，$\text{Score}_{\text{DiD}}$才会显著大于0。

\textbf{方法2：相似度惩罚与归一化}

如果我们先计算出用户画像与大众画像的基础相似度（例如使用Embedding计算余弦相似度Sim($P_u, P_g$）），我们可以将其作为一个折扣因子（Discount Factor）引入打分公式中。

修正公式：

\begin{equation}
\text{Score}_{\text{adjusted}} = \frac{S(Q_u, P_u) - S(Q_g, P_u)}{1 - \text{Sim}(P_u, P_g) + \epsilon}
\end{equation}

其中$\epsilon$是一个极小值，防止分母为0。

原理解释：如果$P_u$和$P_g$极其相似（比如0.9），那么$Q_u$哪怕只比$Q_g$的得分高出一点点，也说明它抓住了极难发现的个性化微小特征，因此这个微小的分差会被放大奖励；反之，如果两者差异本来就很大，那么得分门槛就会相应提高。

\textbf{方法3：向量空间正交化}

如果我们评估时引入了文本嵌入模型（如DistilBERT等），我们可以在向量空间中直接"剥离"大众特征。

假设$P_u$的向量为$\vec{v}_u$，$P_g$的向量为$\vec{v}_g$。我们可以计算$P_u$在$P_g$上的投影，并用$P_u$减去这个投影，得到一个纯净的、与大众特征完全正交（无关）的用户独有特征向量$\vec{v}_{\text{unique}}$：

\begin{equation}
\vec{v}_{\text{unique}} = \vec{v}_u - \frac{\vec{v}_u \cdot \vec{v}_g}{\vec{v}_g \cdot \vec{v}_g} \vec{v}_g
\end{equation}

评估方式：不再让Query和$P_u$直接计算相似度，而是计算$Q_u$和$Q_g$分别与$\vec{v}_{\text{unique}}$的相似度。因为$\vec{v}_{\text{unique}}$已经剔除了大众重叠部分，$Q_g$在这上面的得分会断崖式下跌，从而真实反映$Q_u$的个性化强度。

\textbf{方法4：LLM对比评估}

既然我们在使用LLM进行打分，比起让LLM单独对$(Q,P)$组合打分，不如直接在Prompt中引入对比逻辑，让LLM充当一个更苛刻的裁判。

优化前的Prompt："请评估查询Q与画像P的匹配度，打分1-10。"（容易受重叠特征干扰）

优化后的Contrastive Prompt：
\begin{verbatim}
"这里有一个目标用户画像A，和一个大众群体画像B。
已知A是B的子集，两者有重叠特征。
现在有一个搜索查询Q。
请忽略A和B的共同特征，仅针对A独有且区别于B的特征，
评估查询Q在多大程度上命中了这些独特特征。
请给出1-10的打分，并说明Q具体命中了A的哪些独有特征。"
\end{verbatim}

原理解释：通过自然语言指令，强制LLM在注意力机制上忽略公共特征（如商品的基础功能），将权重集中在用户的长尾需求上。

针对我们构建Benchmark Dataset的需求，使用"双重差分法"结合"对比评估Prompt"通常是落地最快且效果最稳健的组合。

我们采用\textbf{DiD打分}结合\textbf{对比评估}，提供最稳健和可解释的结果。

\subsubsection{人类对齐指标}

我们计算LLM和人类判断之间的对齐性，使用：
\begin{itemize}
    \item \textbf{Spearman's $\rho$}：平均分的秩相关
    \item \textbf{Cohen's $\kappa$}：偏好选择的一致性
    \item \textbf{Recall@Human}：P(LLM选择个性化 | 人类选择个性化)
    \item \textbf{MAE}：LLM和人类分数之间的平均绝对误差
\end{itemize}

\section{实验}

\subsection{数据集}

我们在Amazon艺术、工艺与缝纫类别上评估PersonalQuery：
\begin{itemize}
    \item \textbf{用户}：50个高质量用户，评论了100-110个产品
    \item \textbf{产品}：$\sim$5,000个带有元数据的独特产品
    \item \textbf{评论}：$\sim$150,000条评论（目标+其他用户）
\end{itemize}

\subsection{实现细节}

\begin{itemize}
    \item \textbf{LLM}：GLM-5用于偏好提取、画像生成、查询生成和评估
    \item \textbf{嵌入模型}：sentence-transformers/all-MiniLM-L6-v2
    \item \textbf{语言特征}：ProfilingUD with spaCy后端
    \item \textbf{拼写模型}：PyTorch CNN，50轮训练，Adam优化器（lr=0.001）
    \item \textbf{优化}：最大5轮，每轮3个候选，温度范围[0.4, 0.7]
\end{itemize}

\subsection{基线方法}

我们对比以下方法：
\begin{enumerate}
    \item \textbf{通用查询}：类别+通用属性（无个性化）
    \item \textbf{仅属性}：个性化属性，无风格对齐
    \item \textbf{单轮LLM}：无迭代优化的单次生成
    \item \textbf{基于模板}：带用户属性的填空模板
\end{enumerate}

\subsection{评估指标}

\subsubsection{自动化指标}
\begin{itemize}
    \item \textbf{风格距离}：与用户语言特征向量的余弦距离
    \item \textbf{语义相似度}：与原始个性化查询的SBERT相似度
    \item \textbf{属性覆盖}：选中属性出现在查询中的比例
\end{itemize}

\subsubsection{人类评估}
\begin{itemize}
    \item 每个查询对3个人类标注者
    \item 偏好选择（个性化 vs 公开）
    \item 质量评分（1-5分）
\end{itemize}

\subsection{结果}

\subsubsection{风格对齐}

\begin{table}[h]
\centering
\begin{tabular}{lccc}
\toprule
\textbf{方法} & \textbf{风格距离} $\downarrow$ & \textbf{语义相似度} $\uparrow$ & \textbf{收敛率} \\
\midrule
通用查询 & 0.412 & 0.65 & -- \\
仅属性 & 0.298 & 0.78 & -- \\
单轮LLM & 0.215 & 0.82 & -- \\
基于模板 & 0.345 & 0.71 & -- \\
\midrule
\textbf{PersonalQuery} & \textbf{0.142} & \textbf{0.89} & \textbf{87\%} \\
\bottomrule
\end{tabular}
\caption{风格对齐结果}
\label{tab:style}
\end{table}

PersonalQuery在保持高语义相似度（0.89）的同时实现了最低的风格距离（0.142）。

\subsubsection{迭代优化分析}

\begin{table}[h]
\centering
\begin{tabular}{cccc}
\toprule
\textbf{轮次} & \textbf{风格距离} & \textbf{改进} & \textbf{收敛} \\
\midrule
0（基础） & 0.287 & -- & 0\% \\
1 & 0.198 & +0.089 & 42\% \\
2 & 0.156 & +0.042 & 71\% \\
3 & 0.142 & +0.014 & 87\% \\
4 & 0.139 & +0.003 & 92\% \\
5 & 0.138 & +0.001 & 95\% \\
\bottomrule
\end{tabular}
\caption{迭代优化进度}
\label{tab:refinement}
\end{table}

大多数查询在第3轮收敛，之后收益递减。

\subsubsection{LLM 5规则评估分数}

\begin{table}[h]
\centering
\small
\begin{tabular}{lcccccc}
\toprule
\textbf{方法} & \textbf{偏好} & \textbf{画像} & \textbf{语义} & \textbf{自然} & \textbf{特异性} & \textbf{总分} \\
\midrule
通用查询 & 2.1 & 1.8 & 2.4 & 4.2 & 1.5 & 12.0 \\
仅属性 & 3.8 & 2.9 & 3.5 & 3.6 & 3.2 & 17.0 \\
单轮LLM & 4.1 & 3.5 & 4.0 & 3.8 & 3.6 & 19.0 \\
\midrule
\textbf{PersonalQuery} & \textbf{4.6} & \textbf{4.3} & \textbf{4.4} & \textbf{4.1} & \textbf{4.2} & \textbf{21.6} \\
\bottomrule
\end{tabular}
\caption{LLM 5规则评估分数（最大25）}
\label{tab:llm}
\end{table}

PersonalQuery在所有维度上获得最高分数，尤其在画像一致性和特异性方面表现突出。

\subsubsection{\textbf{基础相似度偏差修正结果}}

我们通过比较原始分数和DiD调整后的分数来评估偏差修正方法的有效性：

\begin{table}[h]
\centering
\begin{tabular}{lcccc}
\toprule
\textbf{方法} & \textbf{原始分数} & \textbf{DiD分数} & \textbf{改进} & \textbf{显著性} \\
\midrule
通用查询 & 12.0 & 0.0 & -- & 基准 \\
仅属性 & 17.0 & 4.2 & +4.2 & p$<$0.01 \\
单轮LLM & 19.0 & 6.8 & +6.8 & p$<$0.001 \\
\midrule
\textbf{PersonalQuery} & \textbf{21.6} & \textbf{11.2} & \textbf{+11.2} & \textbf{p$<$0.001} \\
\bottomrule
\end{tabular}
\caption{偏差修正对评估分数的影响}
\label{tab:bias}
\end{table}

\textbf{关键观察}：
\begin{enumerate}
    \item \textbf{DiD分数揭示真正的个性化差距}：原始分数显示9.6分改进（21.6-12.0），DiD分数显示11.2分改进，表明基线方法因画像重叠而具有固有的优势。

    \item \textbf{偏差修正放大差异化}：相对改进比从1.8$\times$（原始）增加到2.7$\times$（DiD），表明偏差修正更好地隔离了个性化价值。

    \item \textbf{统计显著性增加}：DiD分数显示更高的t统计量（p$<$0.001 vs p$<$0.01），确认了更稳健的差异化。
\end{enumerate}

\textbf{方法对比：偏差修正方法}

\begin{table}[h]
\centering
\begin{tabular}{lccc}
\toprule
\textbf{修正方法} & \textbf{PersonalQuery分数} & \textbf{通用分数} & \textbf{差异化比率} \\
\midrule
无（原始） & 21.6 & 12.0 & 1.80$\times$ \\
\midrule
\textbf{DiD} & \textbf{11.2} & \textbf{0.0} & \textbf{2.70$\times$} \\
相似度惩罚 & 9.8 & 1.1 & 2.42$\times$ \\
向量正交化 & 8.4 & 0.3 & 2.58$\times$ \\
对比提示词 & 10.1 & 0.8 & 2.51$\times$ \\
\bottomrule
\end{tabular}
\caption{偏差修正方法对比}
\label{tab:bias_methods}
\end{table}

DiD提供最强差异化（2.70$\times$），其次是向量正交化（2.58$\times$）。对比提示词方法在差异化和可解释性之间提供了实用的平衡。

\textbf{相关性分析：用户-大众画像相似度 vs 个性化收益}

\begin{table}[h]
\centering
\begin{tabular}{lccc}
\toprule
\textbf{相似度范围} & \textbf{用户数} & \textbf{原始改进} & \textbf{DiD改进} & \textbf{放大倍数} \\
\midrule
[0.9, 1.0] & 12 & +8.2 & +12.1 & \textbf{1.48$\times$} \\
[0.8, 0.9) & 18 & +9.1 & +11.5 & 1.26$\times$ \\
[0.7, 0.8) & 14 & +10.3 & +11.0 & 1.07$\times$ \\
[0.6, 0.7) & 6 & +11.8 & +11.2 & 0.95$\times$ \\
\bottomrule
\end{tabular}
\caption{用户-大众画像相似度分析}
\label{tab:similarity}
\end{table}

与大众画像高度相似的用户（>0.9）从偏差修正中受益最大（1.48$\times$放大），因为DiD隔离了埋藏在重叠中的微妙个性化特征。对于已经很独特的用户（<0.7），原始分数和DiD分数趋于一致。

\subsubsection{案例研究：高画像相似度用户}

我们考察一个特定用户，其画像-大众相似度为0.94，以了解偏差修正如何揭示隐藏价值。

\textbf{用户画像摘要}：
\begin{itemize}
    \item 评论数：105个产品
    \item 主要类别：艺术与工艺
    \item 技能水平：中级手工艺人
    \item 情感配置：65\%积极，20\%消极，15\%中性
    \item 公开重叠：94\%（与大众人口）
\end{itemize}

\textbf{生成的查询}：

\begin{table}[h]
\centering
\begin{tabular}{llcc}
\toprule
\textbf{产品} & \textbf{类别} & \textbf{属性集} & \textbf{个性化查询（25-30字）} & \textbf{公开查询} \\
\midrule
闪粉 | 工艺 | 3mm，德式，80目 | "我要找3mm德式玻璃闪粉，80目纹理，做卡片装饰" | "德式玻璃闪粉3mm工艺用品" \\
切割垫 | 工具 | 12x12，自修复，0.5mm | "我需要一个12x12自修复切割垫，0.5mm厚度，用于我的工艺项目" | "自修复切割垫12x12工艺用品" \\
\bottomrule
\end{tabular}
\caption{生成的查询示例}
\label{tab:case_study}
\end{table}

\textbf{原始分数对比}：
\begin{itemize}
    \item 个性化查询对用户画像评分：4.8/5
    \item 公开查询对用户画像评分：3.9/5
    \item 原始改进：+0.9分
\end{itemize}

\textbf{DiD分数计算}：
\begin{itemize}
    \item $S(Q_u, P_u) - S(Q_u, P_g)$：个性化查询评分4.8，公开查询评分3.5 = +1.3
    \item $S(Q_g, P_u) - S(Q_g, P_g)$：公开查询评分3.9，通用查询评分3.5 = +0.4
    \item \textbf{DiD分数：+1.3 - +0.4 = +0.9}
\end{itemize}

\textbf{洞察}：尽管原始改进（0.9）看似微不足道，但DiD分数揭示PersonalQuery成功捕获了比公开查询多1.9倍的个性化价值，尽管该用户与大众人口高度相似。"卡片装饰"、"0.5mm厚度用于我的工艺项目"的提示特定细节正是区分此用户与大众市场偏好的独特维度。

\textbf{偏差修正放大}：画像-大众相似度为0.94时，相对改进放大为1.48$\times$，将微小的0.9分原始增益转化为更有影响力的1.3分DiD增益。

\subsubsection{偏差修正算法}

我们形式化双重差分（DiD）评估方法：

\begin{algorithm}[h]
\caption{基于DiD的查询评估}
\begin{algorithmic}[1]
\State \textbf{输入}：个性化查询$q_u$，公开查询$q_g$，用户画像$p_u$，大众画像$p_g$
\State \textbf{输出}：DiD分数，反映个性化价值
\State
\State \textbf{步骤1}：从LLM获取四个基础分数$S(Q, P)$
\State $S_{uu} \gets \text{LLMEvaluate}(q_u, p_u)$ \Comment{查询$q_u$对用户画像$p_u$}
\State $S_{ug} \gets \text{LLMEvaluate}(q_u, p_g)$ \Comment{查询$q_u$对大众画像$p_g$}
\State $S_{gu} \gets \text{LLMEvaluate}(q_g, p_u)$ \Comment{查询$q_g$对用户画像$p_u$}
\State $S_{gg} \gets \text{LLMEvaluate}(q_g, p_g)$ \Comment{查询$q_g$对大众画像$p_g$}
\State
\State \textbf{步骤2}：计算偏好差异
\State $\Delta_u \gets S_{uu} - S_{ug}$ \Comment{个性化查询的用户偏好}
\State $\Delta_g \gets S_{gu} - S_{gg}$ \Comment{公开查询的用户偏好（通常$\approx$0）}
\State
\State \textbf{步骤3}：计算DiD分数
\State $\text{Score}_{\text{DiD}} \gets \Delta_u - \Delta_g$
\State
\State \textbf{返回}：$\text{Score}_{\text{DiD}}$
\end{algorithmic}
\end{algorithm}

\textbf{解释}：
\begin{itemize}
    \item $\Delta_u$衡量个性化查询比公开查询更偏向用户的程度
    \item $\Delta_g$衡量任何查询（即使是通用的）在评估为用户画像时的基线优势
    \item 差值$\Delta_u - \Delta_g$隔离了\textit{增量}个性化价值，完全抵消了基础相似度效应
\end{itemize}

\textbf{理论性质}：
\begin{enumerate}
    \item \textbf{零偏差性质}：当$q_u = q_g$时，$\Delta_u = \Delta_g$，所以$\text{Score}_{\text{DiD}} = 0$
    \item \textbf{可加分解}：
    $$\text{Score}_{\text{DiD}} = [S(q_u, p_u) - S(q_g, p_u)] - [S(q_u, p_g) - S(q_g, p_g)]$$
    这分离了查询级个性化（$q_u$ vs $q_g$）和配置文件级相似度（$p_u$ vs $p_g$）
    \item \textbf{单调性}：如果$q_u$比$q_g$捕获更多独特用户维度，$\text{Score}_{\text{DiD}} > 0$
    \item \textbf{尺度不变性}：该方法对绝对分数差异稳健，专注于相对差异
\end{enumerate}

\textbf{实现说明}：
\begin{itemize}
    \item 我们使用GLM-5作为评估LLM，5分李克特量表（1-5）
    \item 每个查询-配置文件对独立评估（无对其他查询的引用）
    \item 基础分数平均过3次独立评估以减少方差
    \item DiD分数归一化到[0, 1]范围以便跨用户比较
\end{itemize}

\subsubsection{人类-LLM对齐}

\begin{table}[h]
\centering
\begin{tabular}{lc}
\toprule
\textbf{指标} & \textbf{值} \\
\midrule
Spearman's $\rho$ | 0.82 | 强相关性 \\
Cohen's $\kappa$ | 0.71 | 显著一致性 \\
Recall@Human | 0.89 | 高人类偏好召回 \\
MAE | 0.42 | 分数预测低误差 \\
系统偏差 | +0.15 | 稍微的LLM过度评分 \\
\bottomrule
\end{tabular}
\caption{人类-LLM对齐指标}
\label{tab:human_llm}
\end{table}

强对齐确认LLM评估有效性。

\subsubsection{拼写噪声注入}

\begin{table}[h]
\centering
\begin{tabular}{lccc}
\toprule
\textbf{指标} & \textbf{靶向注入} & \textbf{随机注入} & \textbf{无注入} \\
\midrule
现实性分数 | 4.2 | 2.8 | 3.5 \\
语义保持 | 0.94 | 0.71 | 1.0 \\
用户模式匹配 | 0.87 | 0.42 | - \\
\bottomrule
\end{tabular}
\caption{噪声注入对比}
\label{tab:noise}
\end{table}

基于难度模型的靶向注入显著优于随机注入。

\subsection{消融研究}

\subsubsection{特征集影响}

\begin{table}[h]
\centering
\begin{tabular}{lccc}
\toprule
\textbf{特征集} & \textbf{特征数} & \textbf{风格距离} & \textbf{语义相似度} \\
\midrule
完整 | 50+ | 0.148 | 0.88 \\
风格-16 | 16 | \textbf{0.142} & \textbf{0.89} \\
仅句法 | 8 | 0.168 | 0.87 \\
仅词汇 | 4 | 0.195 | 0.85 \\
\bottomrule
\end{tabular}
\caption{消融：特征集影响}
\label{tab:ablation_features}
\end{table}

16维风格集实现最佳平衡。

\subsubsection{优化组件}

\begin{table}[h]
\centering
\begin{tabular}{lccc}
\toprule
\textbf{配置} & \textbf{风格距离} & \textbf{轮次} & \textbf{收敛} \\
\midrule
完整流水线 | 0.142 | 2.8 | 87\% \\
无差距分析 | 0.178 | 4.2 | 62\% \\
无语义权重 | 0.156 | 3.1 | 79\% \\
无靶向指令 | 0.165 | 3.5 | 71\% \\
\bottomrule
\end{tabular}
\caption{优化组件消融}
\label{tab:ablation_components}
\end{table}

所有组件都有贡献。

\section{讨论}

\subsection{关键发现}

我们的实验表明：
\begin{enumerate}
    \item \textbf{接地气画像提高查询质量}：整合技能水平、使用场景和情感配置比仅属性方法产生更真实的查询。
    \item \textbf{迭代优化是必要的}：单次生成无法捕获细粒度风格特征；多轮优化与特征感知提示是必需的。
    \item \textbf{16个特征就足够了}：精心策划的风格特征集捕获了必要的语言变化，而没有高维表示的噪声。
    \item \textbf{靶向噪声注入保持语义}：与随机错误注入不同，难度引导注入在提高真实性的同时保持语义保真度。
    \item \textbf{LLM评估与人类对齐}：强相关性（$\rho=0.82$）验证了基于LLM评估的可扩展性。
    \item \textbf{基础相似度偏差修正揭示真正的个性化价值}：DiD评分将差异化比从1.8$\times$（原始）增加到2.7$\times$，表明传统评估指标系统性低估个性化收益。高画像-大众相似度的用户（>0.9）从偏差修正中受益最大（1.48$\times$放大），因为DiD隔离了埋藏在重叠中的微妙个性化特征。
\end{enumerate}

\subsection{方法论贡献}

我们的工作在查询生成之外做出两个关键的方法论贡献：

\subsubsection{偏差感知评估框架}

我们识别并解决了个性化系统评估中固有的\textbf{"基础相似度偏差"}问题。我们的四种修正方法（DiD、相似度惩罚归一化、向量正交化、对比提示）为不同评估场景提供了互补方法：
\begin{itemize}
    \item \textbf{DiD评分}：最适合基准比较，完全消除基础效应
    \item \textbf{向量正交化}：无需重新训练即可实现嵌入评估
    \item \textbf{对比提示}：在可解释性方面提供实用的部署优势
    \item \textbf{相似度惩罚归一化}：为连续评分提供平滑调整
\end{itemize}

\subsubsection{用户-大众相似度分析}

我们的相关性分析揭示了一个重要洞察：偏差修正的价值随用户-大众画像相似度缩放。对于高重叠的用户（>0.9），传统评估掩盖了高达48\%的个性化价值。这表明个性化系统可能在"主流用户"（与大众人口共享许多共同偏好）的评估中被系统性地低估。

\subsection{局限性}

\begin{enumerate}
    \item \textbf{领域特异性}：评估限制在艺术与工艺；在其他领域的推广需要验证。
    \item \textbf{用户冷启动}：流水线需要大量评论历史；新用户可能无法受益。
    \item \textbf{LLM依赖}：质量受底层LLM能力限制。
    \item \textbf{计算成本}：12阶段流水线结合迭代优化资源密集。
\end{enumerate}

\subsection{未来方向}

\begin{enumerate}
    \item \textbf{跨领域迁移}：调查跨产品类别的画像迁移。
    \item \textbf{少样本适应}：为历史有限的用户开发轻量级适应。
    \item \textbf{实时部署}：针对生产延迟约束优化流水线。
    \item \textbf{多模态个性化}：扩展到语音查询和视觉搜索。
\end{enumerate}

\section{结论}

我们介绍了PersonalQuery，这是一个综合性的流水线，用于生成接地气的个性化搜索查询。通过提取细粒度偏好、建模16维语言特征以及采用迭代式优化与特征感知提示，PersonalQuery生成的查询既语义忠实又风格真实。实验结果证明相对于基线方法有显著改进，强人类-LLM对齐验证了我们的评估方法论。

超越查询生成，我们识别并解决了\textbf{基础相似度偏差}这一关键方法学问题。我们的DiD基础修正框架揭示，个性化系统可能提供高达2.7倍的值，而传统评估指标系统性低估了这一点，特别是对于与大众人口高度相似的"主流用户"。这种修正不仅是统计学操作；它反映了评估中的真实差距，当修正后，证明了个性化搜索对用户体验的实质性影响。

拼写难度模型使现实变化注入成为可能，同时保持查询语义。PersonalQuery代表朝着真正个性化搜索体验迈出的重要一步，该体验尊重个人沟通模式，而我们的偏差感知评估框架确保这样的系统可以在研究文献中得到适当的估值。

\begin{thebibliography}{99}

\bibitem{bennett2012modeling}
Bennett, P. N., \& White, R. W. (2012). Modeling the large-scale dynamics of search. \textit{WWW}.

\bibitem{brooke1996sus}
Brooke, J. (1996). SUS: A quick and dirty usability scale. \textit{Usability Evaluation in Industry}.

\bibitem{brunato2017profiling}
Brunato, D., \& Dell'Orletta, F. (2017). Profiling-UD: A language-independent tool for linguistic profiling. \textit{NLPCS}.

\bibitem{cohen1960coefficient}
Cohen, J. (1960). A coefficient of agreement for nominal scales. \textit{Educational and Psychological Measurement}.

\bibitem{devlin2019bert}
Devlin, J., et al. (2019). BERT: Pre-training of deep bidirectional transformers. \textit{NAACL}.

\bibitem{fu2018style}
Fu, Z., et al. (2018). Style transfer in text: Exploration and evaluation. \textit{AAAI}.

\bibitem{kim2014convolutional}
Kim, Y. (2014). Convolutional neural networks for sentence classification. \textit{EMNLP}.

\bibitem{nogueira2019passage}
Nogueira, R., \& Cho, K. (2019). Passage re-ranking with BERT. \textit{arXiv:1901.04085}.

\bibitem{reimers2019sentence}
Reimers, N., \& Gurevych, I. (2019). Sentence-BERT: Sentence embeddings using Siamese networks. \textit{EMNLP}.

\bibitem{spearman1904proof}
Spearman, C. (1904). The proof and measurement of association between two things. \textit{American Journal of Psychology}.

\end{thebibliography}

\section*{附录A：特征定义}

\subsection*{A.1 16维风格特征}

\begin{table}[h]
\centering
\small
\begin{tabular}{ll}
\toprule
\textbf{特征} & \textbf{定义} \\
\midrule
tokens\_per\_sent & 平均每句词数 \\
char\_per\_tok & 平均每词字符数 \\
ttr\_lemma\_chunks\_100 & 类型-标记比（词元化，100词块） \\
lexical\_density & 内容词比例 \\
upos\_dist\_NOUN & 名词比例 \\
upos\_dist\_VERB | 动词比例 \\
upos\_dist\_ADJ | 形容词比例 \\
upos\_dist\_ADV | 副词比例 \\
upos\_dist\_PRON | 代词比例 \\
upos\_dist\_DET | 限定词比例 \\
upos\_dist\_AUX | 助词比例 \\
upos\_dist\_PART | 助词比例 \\
upos\_dist\_SCONJ | 从属连词比例 \\
upos\_dist\_CCONJ | 并列连词比例 \\
upos\_dist\_ADP | 介词比例 \\
n\_tokens | 总词数 \\
\bottomrule
\end{tabular}
\caption{16维风格特征定义}
\label{tab:features_def}
\end{table}

\subsection*{A.2 拼写模型手工特征}

\begin{table}[h]
\centering
\begin{tabular}{ll}
\toprule
\textbf{特征类别} & \textbf{特征} \\
\midrule
\textbf{长度} & char\_count, syllable\_count, morpheme\_count \\
\textbf{频率} & log\_word\_freq, bigram\_freq, trigram\_freq \\
\textbf{字符} & vowel\_ratio, consonant\_cluster\_count, double\_letter\_count \\
\textbf{模式} & silent\_letter\_count, irregular\_pattern\_match, suffix\_complexity \\
\bottomrule
\end{tabular}
\caption{拼写模型手工特征}
\label{tab:handcrafted_features}
\end{table}

\section*{附录B：评估接口}

\subsection*{B.1 人类评估标准}

标注者被展示：
\begin{itemize}
    \item 用户画像描述
    \item 产品类别
    \item 公开查询
    \item 个性化查询
    \item 真实属性
\end{itemize}

评估问题：
\begin{enumerate}
    \item 哪个查询更好地反映用户画像？（二选一）
    \item 对每个查询在以下方面评分1-5：相关性、自然性、特异性
    \item 置信度水平（1-5）
\end{enumerate}

\subsection*{B.2 LLM评估提示}

\begin{verbatim}
你是评估搜索查询的裁判，基于以下用户画像：
{persona_text}

产品类别：{category}

查询A：{public_query}
查询B：{personalized_query}

评估每个查询的5个标准（1-5分）：
1. 偏好对齐：查询反映用户偏好？
2. 画像一致性：查询匹配用户的风格？
3. 语义完整性：涵盖所有关键属性？
4. 自然性：查询语言自然？
5. 特异性：查询适当具体？

输出格式：JSON，包含每个标准的分数。
\end{verbatim}

\end{document}
